\medskip\noindent\textbf{Supplementary Note 3. Data preparation and scIB scores.}\par\smallskip
\begingroup\small\setlength{\parskip}{3pt}
The eight tasks come from scIB, figshare article 12420968, version 8 (Supplementary Table S1). All cells are retained. RNA tasks use scIB's batch and cell-type variables and 2,000 highly variable genes selected batch-wise with cell-ranger flavour and 20 bins on normalised expression. Chromatin-accessibility tasks use their supplied labels and all 3,429 or 3,580 gene-activity features. Counts are used as distributed, including non-integer values in 18 batches. These comprise transcripts per million in the Villani batch of both immune tasks, reads per kilobase per million, expected counts and collision-corrected counts in the pancreas batches smarter, fluidigmc1, celseq and celseq2, and normalised values in the twelve 10x batches of lung.

scib 1.1.7 scores posterior means for all cells, including training cells. Metrics retain their raw scales, with higher values better. The batch score $B$ averages five batch metrics. The biology score $C$ averages the applicable biology metrics, and $S=0.4B+0.6C$ is the selection score. Unlike scIB's overall score, these means use no scaling across methods. Let $K$ denote batches, $L$ cell types and $G_{15}$ the embedding's 15-nearest-neighbour graph.

\noindent\textit{Batch metrics.} Principal-component regression is $\mathrm{PCR}=\sum_j R_j^2V_j/\sum_jV_j$, where $R_j^2$ measures batch association and $V_j$ component variance. It uses 50 components of unintegrated normalised expression on all genes and all embedding dimensions. Its score is $\max\{0,(\mathrm{PCR}_{\mathrm{before}}-\mathrm{PCR}_{\mathrm{after}})/\mathrm{PCR}_{\mathrm{before}}\}$.

Batch average silhouette width (ASW) averages $1-|a_i|$ within cell types spanning multiple batches, then across types. Silhouettes use Euclidean distances and batch clusters. Graph iLISI uses 90 shortest-path neighbours on $G_{15}$, weighted at perplexity 30, for a random subsample of about half the cells. With median inverse Simpson index $m$, its score is $(m-1)/(K-1)$. Graph connectivity averages, over cell types, the fraction of a type's cells in the largest connected component of $G_{15}$ restricted to that type.

kBET uses a 50-nearest-neighbour graph and tests types with at least 10 cells spanning multiple batches. Its $\chi^2$ test compares diffusion-neighbour batch composition with the type's composition, using $k_0=\min\{70,\max(10,\lfloor\bar n_c/4\rfloor)\}$, where $\bar n_c$ is mean cells per batch. Disconnected components smaller than $3k_0$ are dropped if at least 75\% of cells remain. Otherwise the rejection rate is 1. The score is one minus mean rejection rate across tested types. R's seed is 0 for each call.

\noindent\textit{Biology metrics.} Leiden clustering of $G_{15}$ uses resolutions $0.2,0.4,\dots,2.0$. The clustering $\kappa$ maximising cell-type normalised mutual information is retained, with $\mathrm{NMI}=2I(\kappa;c)/(H(\kappa)+H(c))$. Its adjusted Rand index is $\big(\sum_{uv}\binom{n_{uv}}2-\Sigma_{ab}\big)/\big(\tfrac12[\sum_u\binom{a_u}2+\sum_v\binom{b_v}2]-\Sigma_{ab}\big)$, where $n_{uv}$ counts cluster--type intersections, $a_u,b_v$ are margins and $\Sigma_{ab}=\sum_u\binom{a_u}2\sum_v\binom{b_v}2/\binom n2$.

Cell-type ASW is $(\bar a+1)/2$, using mean Euclidean silhouette by cell type. Isolated labels occur in the fewest batches, unless every type occurs in every batch. Isolated-label F1 averages each label's best F1 over clusters and resolutions. Isolated-label ASW averages $(\bar a_\ell+1)/2$ over those labels, where $\bar a_\ell$ is the mean cell-type silhouette within a label. Graph cLISI uses iLISI neighbourhoods with cell-type labels and score $(L-m)/(L-1)$.

Cell-cycle conservation applies to pancreas, lung and both immune tasks. Human S and G2M gene scores give PCR variance contributions $v$ before and after integration within each batch. The score averages $\max\{0,1-|v_{\mathrm{after}}-v_{\mathrm{before}}|/v_{\mathrm{before}}\}$ across batches. Trajectory conservation applies to the immune tasks. Diffusion pseudotime is recomputed within the largest connected component, using a root chosen from supplied pseudotime and labels. Values above 1 and outside the component become 0. The score is $(\rho+1)/2$ for Spearman correlation with supplied pseudotime, or 0 without a root.

All tasks use the six metrics from NMI through cLISI. Pancreas and lung add cell-cycle conservation, giving seven metrics. Both immune tasks also add trajectory conservation, giving eight. Undefined metrics are omitted. Highly variable gene conservation is undefined for embeddings. All batch metrics are available. Scores lie between 0 and 1 except ARI, whose minimum was $-0.00003$.
\par\endgroup
